\section{Cuts adopted and inputs for the mixed-event step}
\label{sec:cuts}

\begin{table}[h]
\centering
\caption{Selection used in this QA (defaults in \fpath{macros/runTofMEQA.C} and the constructor of \fcmd{StTofMEQAMaker}).}
\label{tab:cuts}
\small
\begin{tabular}{llp{5cm}}
\toprule
Level & Cut & Status \\
\midrule
Event & trigger $\in$ \{640001, 640011, 640021, 640031, 640041, 640051\} & applied \\
Event & $|V_z|<70$\,cm & applied \\
Event & $V_r<2$\,cm & applied \\
Event & $|V_z-V_z^{\rm VPD}|$ & not applied \\
Event & pile-up / bad-run rejection & not implemented \\
Track & primary, $p_T>0.15$\,\GeVc & applied \\
Track & nHitsFit $\ge15$, nHitsFit/nHitsMax $>0.52$ & applied \\
Track & gDCA $<3$\,cm & applied \\
Track & $|\eta|<1$ & applied \\
TOF & btofMatchFlag $>0$ & applied \\
TOF & $\beta>0$ & applied (TOF-matched sample) \\
TOF & $|y_{\rm local}|$, $|z_{\rm local}|$ & none beyond production ($|z_{\rm local}|<3.05$\,cm) \\
\bottomrule
\end{tabular}
\end{table}

What this QA implies for the mixed-event matching:
\begin{itemize}
\item With \hitsPerMatched{} stored hits per TOF-matched QA track and a matching window of one pad in $y$
  and $\pm3.05$\,cm in $z$, the random-match probability will depend strongly on event multiplicity;
  events must be mixed within multiplicity classes (refMult or number of BTOF hits) and $V_z$ classes.
\item The picoDst stores only the cell ID of each BTOF hit, so matching is done at the cell level with the
  projection of Sec.~\ref{sec:me} (implemented).
\item Dead trays/modules are identical between mixed events because mixing is restricted to one run (implemented).
\item The $y_{\rm local}$ side-bands ($|y_{\rm local}|>2.5$\,cm) provide an independent same-event
  estimate of the random-match level that the mixed-event result can be compared to.
\end{itemize}

\decision{(1) Is the minimum-bias trigger list correct for P24iy, or should the analysis use 640051 only?
(2) Should a $|V_z-V_z^{\rm VPD}|$ cut be used, given the VPD vertex is available for only \fracVpdValid\% of
events? (3) Mixing classes: refMult (used now) or number of BTOF hits, and is a 10\,cm $V_z$ class adequate?
(4) Which matching rule should the final random-match correction use: production-like cell$\pm$1 (default
proposal) or cell only?
(5) Low-$p_T$ projection: accept the DCA-helix projection, or add an energy-loss correction / use the
production cell for same-event matches?}
